\documentclass[11pt,reqno]{amsart}

\usepackage[T1]{fontenc}
\usepackage[utf8]{inputenc}
\usepackage{amsmath,amssymb,amsthm}
\usepackage{mathtools}
\usepackage{xcolor}
\usepackage[margin=1.15in]{geometry}
\usepackage[colorlinks=true,linkcolor=blue!55!black,citecolor=blue!55!black,
            urlcolor=blue!55!black]{hyperref}
\usepackage[nameinlink,capitalise]{cleveref}
\usepackage[expansion=false]{microtype}

\theoremstyle{plain}
\newtheorem{theorem}{Theorem}[section]
\newtheorem{proposition}[theorem]{Proposition}
\newtheorem{lemma}[theorem]{Lemma}
\newtheorem{corollary}[theorem]{Corollary}

\theoremstyle{definition}
\newtheorem{definition}[theorem]{Definition}
\newtheorem{example}[theorem]{Example}

\theoremstyle{remark}
\newtheorem{remark}[theorem]{Remark}

\numberwithin{equation}{section}

\DeclareMathOperator{\Image}{Im}
\DeclareMathOperator{\Spec}{Spec}
\DeclareMathOperator{\Real}{Re}
\DeclareMathOperator{\id}{id}

\begin{document}

\title[Canonical Minimal Source Completion]{Canonical Minimal Source Completion:\\
The Coarsest Readout Extension on which a Nominated\\
Family of Source Relations Becomes Well Defined}

\author{Jeremy Rodgers}
\thanks{Website: \href{https://everythingequation.com}{everythingequation.com}}
\address{Independent Researcher}

\date{June 2026}

\subjclass[2020]{Primary 18A30, 54B15; Secondary 53C29, 58J50, 81Q70}

\keywords{Quotient construction, universal property, terminal object, relation-sufficient
extension, quotient topology, gauge reduction, holonomy, flat connections, spectral
invariants, descent}

\begin{abstract}
Let $p\colon S\to T$ be a surjective source-to-readout map on a reduced source space, and let
$\{r_\alpha\colon S\to V_\alpha\}_{\alpha\in A}$ be a set-indexed family of physically invariant
source relations that need not descend along $p$. We construct the coarsest extension of the
readout on which every $r_\alpha$ becomes well defined. The construction is explicit: form the
joint map $j_A=(p,(r_\alpha)_{\alpha\in A})\colon S\to T\times\prod_{\alpha\in A}V_\alpha$ and
set $E_A=\Image j_A$; equivalently, quotient $S$ by equality of the readout together with every
selected relation value. We prove that $E_A$ is terminal in the category of surjective
relation-sufficient extensions: every sufficient extension $E$ admits a unique surjection
$E\to E_A$. This yields existence, uniqueness up to unique isomorphism, and a precise minimality
theorem, without nominating a preferred representative of any readout fiber. We characterize
redundant and indispensable relation coordinates by an exact separation criterion, construct
canonical refinement maps for nested families and prove their cocycle law, prove invariance under
faithful reparameterization of value spaces, and construct the completion determined directly by a
family of target-answer maps, together with the canonical surjection from the relation completion
onto it. We give the gauge-reduced form of the theorem (completion cannot resurrect redundancy
already quotiented away) and the topological form, distinguishing carefully between the quotient
topology on $E_A$ and the subspace topology inherited from the product. For flat $U(1)$ connections
on a circle, whose curvature readout is identically zero, oriented transport completion returns
$U(1)$, whereas the minimal completion for the charged-scalar spectral target is the inversion
quotient $U(1)/(z\sim z^{-1})\cong[-1,1]$; the canonical refinement $e^{i\theta}\mapsto\cos\theta$
exhibits the general relation-to-target surjection concretely.
\end{abstract}

\maketitle

\section{Introduction}
\label{sec:intro}

\subsection{The problem inherited from Papers 1 and 2}

This is the third paper of a sequence on source--readout non-equivalence. The first companion
manuscript \cite{Paper1} established that a readout map $p\colon S\to T$ need not admit descent for
a given source-defined quantity: a function $r\colon S\to V$ factors as $r=\bar r\circ p$ if and
only if $r$ is constant on every fiber $F_t=p^{-1}(t)$, and equivariant refinements of this
statement produce genuine obstructions. The second companion manuscript \cite{Paper2} sharpened the
significance of that failure. Information destroyed by a readout is not by itself a defect; it
becomes consequential only \emph{relative to a nominated target}. Paper~2 characterized exactly
when a target with correct-answer map $a\colon S\to A_Q$ is solvable from the readout alone
(precisely when $a$ is $\sim_p$-invariant), and, when it is not, constructed the coarsest extension
of $T$ sufficient to determine that single answer map.

What Paper~2 does not supply is the family-level construction. Physical practice almost never
nominates a single answer map. One nominates a \emph{family}: a collection of invariant source
relations $r_\alpha$, indexed by a set $A$, each entering some independently specified mathematical
or physical target. Assembling the single-target completions of Paper~2 one at a time does not
answer the family question, because the completions do not simply compose: the object one wants is
a single extension of the readout on which \emph{all} the nominated relations become simultaneously
well defined, and one wants it to be canonical---adding no distinctions beyond those the nominated
family forces. Neither existence, nor minimality, nor uniqueness of such an object is formal from
Paper~2.

\subsection{What this paper proves}

Fix the reduced source $S$, the surjective readout $p\colon S\to T$, and a set-indexed family
$r_\alpha\colon S\to V_\alpha$, $\alpha\in A$. Write $V_A=\prod_{\alpha\in A}V_\alpha$,
$R_A=(r_\alpha)_{\alpha\in A}\colon S\to V_A$, and
\begin{equation}
\label{eq:jA-intro}
j_A=(p,R_A)\colon S\longrightarrow T\times V_A,
\qquad
E_A:=\Image(j_A).
\end{equation}
Let $e_A\colon S\twoheadrightarrow E_A$ be the induced surjection and let
$u_A\colon E_A\to T$, $\rho_\alpha\colon E_A\to V_\alpha$ be the coordinate projections. Our
results are as follows.

\begin{enumerate}
\item \emph{Explicit construction and quotient realization} (\cref{sec:construction}). We have
$p=u_A\circ e_A$ and $r_\alpha=\rho_\alpha\circ e_A$ for all $\alpha$, and $E_A\cong S/\!\sim_A$
where $s\sim_A s'$ iff $p(s)=p(s')$ and $r_\alpha(s)=r_\alpha(s')$ for every $\alpha\in A$. The
information ordering $\Delta_S\subseteq\ \sim_A\ \subseteq\ \sim_p$ holds, and the fiber
$u_A^{-1}(t)$ is canonically the set $R_A(F_t)$ of \emph{realized relation signatures} over $t$.

\item \emph{Terminality} (\cref{sec:category,sec:main}). We define the category of surjective
relation-sufficient extensions and prove that $\mathcal{E}_A$ is \emph{terminal} in it: for every
sufficient extension $\mathcal E=(E,e,u,(\rho_\alpha^E))$ there is a \emph{unique} morphism
$\phi_{\mathcal E}\colon E\to E_A$, and it is surjective. The arrow runs $E\to E_A$, never
$E_A\to E$. Elementarily, terminality is the statement $\sim_e\ \subseteq\ \sim_A$ for every
sufficient extension: $E_A$ is the \emph{coarsest} sufficient extension, not the finest.

\item \emph{Coordinate economy} (\cref{sec:redundancy}). A nominated relation $r_\alpha$ is
indispensable if and only if some pair $s,s'$ agrees on the readout and on all
$r_\beta$, $\beta\neq\alpha$, yet is separated by $r_\alpha$; and deleting $r_\alpha$ strictly
coarsens the completion exactly in that case.

\item \emph{Functoriality in the family} (\cref{sec:refinement}). For $A\subseteq B$ there is a
unique surjection $\pi_{BA}\colon E_B\to E_A$ over $S$, and these satisfy
$\pi_{CA}=\pi_{BA}\circ\pi_{CB}$. Faithful reparameterization of value spaces leaves the completion
canonically unchanged over $T$.

\item \emph{Target completion} (\cref{sec:targets}). The completion $E_{\mathcal Q}$ determined by a
family of answer maps is constructed, and when the relations mediate the answers there is a
canonical surjection $\chi\colon E_A\to E_{\mathcal Q}$, an isomorphism precisely when the two
induced equivalence relations coincide. Target completion may be \emph{strictly coarser} than
relation completion.

\item \emph{Gauge boundary and topology} (\cref{sec:gauge,sec:topology}). Completion occurs
\emph{after} gauge reduction and cannot undo it. In the topological setting, $E_A$ with the quotient
topology is terminal among topological sufficient extensions; the canonical injection
$E_A\hookrightarrow T\times V_A$ is an embedding exactly when quotient and subspace topologies
agree, for which compactness of $S$ and Hausdorffness of $T\times V_A$ suffice.

\item \emph{Worked example} (\cref{sec:u1}) and \emph{exact reconstruction scope}
(\cref{sec:reconstruction}): $e_A$ is injective iff $\{p\}\cup\{r_\alpha\}$ separates points of $S$.
\end{enumerate}

\subsection{What this paper does not prove}
\label{subsec:scope-front}

The theorem is \emph{model-relative}. Once a reduced source space, a readout, and a nominated
family of invariant relations or targets are specified, there is a canonical coarsest sufficient
quotient. Nothing here shows that the source space is empirically real, that observed physics is a
projection of a hidden source, that every lost distinction must be restored, that the full source is
in general recoverable, or that any particular relation family is physically privileged absent
independent justification. In particular, ``minimal completion'' is not ``reconstruction of the
source'': the two coincide only under the separation criterion of
\cref{prop:reconstruction}. We return to these limits in \cref{sec:discussion}.

\subsection{Notation}

We keep the frozen notation of the companion manuscripts. $S$ is the reduced source space, $T$ the
readout space, $p\colon S\to T$ the surjective readout, and $F_t=p^{-1}(t)$ the readout fiber over
$t\in T$. Throughout, $A$ denotes a \emph{set} indexing the nominated relation family; products
$V_A=\prod_{\alpha\in A}V_\alpha$ are ordinary set-theoretic products, so that a point of $V_A$ is a
choice function $\alpha\mapsto v_\alpha\in V_\alpha$, and no cardinality restriction on $A$ is
needed. We write $\Delta_S=\{(s,s):s\in S\}$ for the equality relation on $S$ and $\sim_p$ for
equality of the readout. For equivalence relations, $\sim\ \subseteq\ \approx$ means that $\sim$ is
\emph{finer} (identifies fewer pairs, retains more information) and $\approx$ is \emph{coarser}.

\section{Reduced source, readout, and the nominated relation family}
\label{sec:setup}

\subsection{The reduced source}

We assume, as in Paper~1, that $S$ has already been reduced: all declared gauge redundancy and all
coordinate-labelling redundancy have been quotiented out. Concretely, one begins with a raw
configuration space $\Omega$ carrying an equivalence relation $\simeq_{\mathrm{phys}}$ of physical
indistinguishability---gauge equivalence, diffeomorphism relabelling, choice of chart---and forms
\begin{equation}
\label{eq:gauge-quotient}
q\colon\Omega\longrightarrow S=\Omega/\!\simeq_{\mathrm{phys}}.
\end{equation}
The order of operations matters and is not a convention we are free to reverse: \emph{reduction
precedes completion}. \Cref{sec:gauge} states the consequence precisely.

\begin{definition}[Eligible relation]
\label{def:eligible}
A raw quantity $\widetilde r\colon\Omega\to V$ is \emph{eligible} for source completion when it is
invariant on physical-equivalence classes,
\begin{equation}
\label{eq:eligible}
\omega\simeq_{\mathrm{phys}}\omega'
\quad\Longrightarrow\quad
\widetilde r(\omega)=\widetilde r(\omega'),
\end{equation}
so that there is a unique $r\colon S\to V$ with $\widetilde r=r\circ q$. All relations nominated
below are assumed eligible, i.e.\ genuinely defined on the reduced source $S$.
\end{definition}

Eligibility is exactly the condition under which the raw quantity is a function of the physical
state rather than of a representative of it. A quantity failing \eqref{eq:eligible}---a gauge
potential, a component in a chosen frame, a coordinate label---is not a candidate relation, and no
completion of the readout can make it one.

\subsection{External nomination}

Fix a set $A$ and a family of eligible relations
\begin{equation}
\label{eq:family}
r_\alpha\colon S\longrightarrow V_\alpha,\qquad \alpha\in A .
\end{equation}
We stress the epistemic status of the family, because the mathematics below is only as meaningful as
the nomination that feeds it. The relations $r_\alpha$ are nominated \emph{externally}: each is
selected because it enters an independently specified mathematical or physical target---an
amplitude, a spectrum, a transport law, an interference pattern. Physical relevance is
\textbf{not} inferred from fiber variation. That a function happens to be nonconstant on $F_t$ shows
only that the readout does not determine it; it does not show that the distinctions it makes are
physical. Confusing the two would convert every artifact of a redundant description into a physical
observable. The theorems below are therefore conditional in exactly one place, and it is the right
place: they take the family as given and tell one, completely, what follows from it.

\subsection{The completion problem}

Write
\begin{equation}
\label{eq:VA-RA}
V_A:=\prod_{\alpha\in A}V_\alpha,
\qquad
R_A\colon S\to V_A,
\qquad
R_A(s):=(r_\alpha(s))_{\alpha\in A}.
\end{equation}

\begin{definition}[The family completion problem]
\label{def:problem}
An extension of the readout \emph{sufficient} for the family \eqref{eq:family} is a space $E$
through which $S$ maps onto, such that
\begin{enumerate}
\item the original readout is retained: $p$ factors through $E$; and
\item every nominated relation becomes a well-defined function of the completed state: each
$r_\alpha$ factors through $E$.
\end{enumerate}
The problem is to construct the \emph{coarsest} such extension: one that makes no distinction
between source states beyond those forced by requirements (1) and (2), and to prove it canonical.
\end{definition}

A sufficient extension is emphatically \emph{not} required to distinguish source states on which
both the readout and every selected relation agree. This is the entire content of minimality, and it
is what separates completion from reconstruction.

\begin{remark}[The empty family]
\label{rem:empty}
If $A=\varnothing$, the ordinary set-theoretic convention gives $V_\varnothing=\{*\}$, the singleton
empty product, and $j_\varnothing=(p,\ast)$, so that
$E_\varnothing=\Image(j_\varnothing)\cong T$ and $u_\varnothing$ is a bijection. Nominating nothing
completes nothing: the completion of the readout relative to the empty family is the readout. This
is the degenerate boundary case of every theorem below, and we will not comment on it again.
\end{remark}

\section{The joint-image construction and its quotient realization}
\label{sec:construction}

\begin{definition}[Minimal relation completion]
\label{def:EA}
Define the \emph{joint map}
\begin{equation}
\label{eq:jA}
j_A\colon S\longrightarrow T\times V_A,
\qquad
j_A(s)=\bigl(p(s),R_A(s)\bigr),
\end{equation}
and define the \emph{minimal relation completion}
\begin{equation}
\label{eq:EA}
E_A:=\Image(j_A)\ \subseteq\ T\times V_A .
\end{equation}
Let $e_A\colon S\twoheadrightarrow E_A$, $e_A(s):=j_A(s)$, be the induced surjection, and let
\begin{equation}
\label{eq:structure-maps}
u_A\colon E_A\longrightarrow T,
\qquad
\rho_\alpha\colon E_A\longrightarrow V_\alpha
\qquad(\alpha\in A)
\end{equation}
be the restrictions to $E_A$ of the coordinate projections of $T\times V_A$.
\end{definition}

\begin{proposition}[Factorization]
\label{prop:factorization}
The maps of \cref{def:EA} satisfy
\begin{equation}
\label{eq:factorization}
p=u_A\circ e_A,
\qquad
r_\alpha=\rho_\alpha\circ e_A
\qquad(\alpha\in A),
\end{equation}
and $e_A$ is surjective. In particular every nominated relation descends to $E_A$.
\end{proposition}

\begin{proof}
Surjectivity of $e_A$ is immediate, since $E_A$ was defined as the image of $j_A$ and $e_A$ is $j_A$
with codomain restricted to that image. For $s\in S$ we have
$e_A(s)=(p(s),(r_\alpha(s))_{\alpha\in A})$. Applying the first coordinate projection gives
$u_A(e_A(s))=p(s)$, and applying the projection onto the $\alpha$-th factor of $V_A$ gives
$\rho_\alpha(e_A(s))=r_\alpha(s)$. As $s$ was arbitrary, \eqref{eq:factorization} holds as an
identity of functions on $S$.
\end{proof}

Thus $(E_A,e_A,u_A,(\rho_\alpha)_{\alpha\in A})$ solves the \emph{existence} half of
\cref{def:problem}. The content of the paper is that it solves the \emph{minimality} and
\emph{uniqueness} halves as well.

\subsection{The induced equivalence relation}

\begin{definition}
\label{def:simA}
Define a relation $\sim_A$ on $S$ by
\begin{equation}
\label{eq:simA}
s\sim_A s'
\quad:\Longleftrightarrow\quad
\bigl[\,p(s)=p(s')\ \text{ and }\ r_\alpha(s)=r_\alpha(s')\ \text{ for every }\alpha\in A\,\bigr].
\end{equation}
\end{definition}

This is an equivalence relation, being an intersection of pullbacks of equality relations along
functions: reflexivity, symmetry, and transitivity for $\sim_A$ follow coordinatewise from the
corresponding properties of equality in $T$ and in each $V_\alpha$.

\begin{proposition}[Quotient realization]
\label{prop:quotient}
The assignment
\begin{equation}
\label{eq:kappa}
\kappa\colon S/\!\sim_A\ \longrightarrow\ E_A,
\qquad
\kappa\bigl([s]_A\bigr):=j_A(s)
\end{equation}
is a well-defined bijection, and it satisfies $\kappa\circ\pi_A=e_A$, where
$\pi_A\colon S\to S/\!\sim_A$ is the canonical projection. Hence
\begin{equation}
\label{eq:EA-quotient}
E_A\ \cong\ S/\!\sim_A
\end{equation}
canonically.
\end{proposition}

\begin{proof}
We verify the three assertions separately, rather than invoking a generic isomorphism
principle for quotients, since the well-definedness step is precisely where the definition
of the relation $\sim_A$ earns its keep.

\emph{Well-definedness.} Suppose $[s]_A=[s']_A$, i.e.\ $s\sim_A s'$. By \eqref{eq:simA} this means
$p(s)=p(s')$ and $r_\alpha(s)=r_\alpha(s')$ for every $\alpha\in A$. Two elements of the product
$T\times V_A$ are equal precisely when all of their coordinates agree; and here the coordinates of
$j_A(s)$ and $j_A(s')$ are $p(s),p(s')$ in the $T$-factor and $r_\alpha(s),r_\alpha(s')$ in the
$\alpha$-th factor of $V_A$, respectively. All agree, so $j_A(s)=j_A(s')$. Hence the value assigned
by \eqref{eq:kappa} does not depend on the chosen representative, and $\kappa$ is a well-defined
function.

\emph{Injectivity.} Suppose $\kappa([s]_A)=\kappa([s']_A)$, i.e.\ $j_A(s)=j_A(s')$. Comparing
coordinates in $T\times V_A$ gives $p(s)=p(s')$ and $r_\alpha(s)=r_\alpha(s')$ for every $\alpha\in
A$, which is exactly \eqref{eq:simA}. Hence $s\sim_A s'$ and $[s]_A=[s']_A$.

\emph{Surjectivity.} Let $x\in E_A$. By \cref{def:EA}, $E_A=\Image(j_A)$, so there is $s\in S$ with
$x=j_A(s)=\kappa([s]_A)$.

Finally, for every $s\in S$ we have $\kappa(\pi_A(s))=\kappa([s]_A)=j_A(s)=e_A(s)$, so
$\kappa\circ\pi_A=e_A$.
\end{proof}

So the joint image and the quotient by ``agreement of the readout and of every nominated relation''
are the same object, viewed extensionally and intensionally. The first description makes existence
manifest; the second makes minimality manifest.

\subsection{The information ordering}

\begin{proposition}[Information ordering]
\label{prop:ordering}
As relations on $S$,
\begin{equation}
\label{eq:ordering}
\Delta_S\ \subseteq\ \sim_A\ \subseteq\ \sim_p ,
\end{equation}
and consequently there is a factorization of the readout,
\begin{equation}
\label{eq:tower}
S\ \xrightarrow{\ e_A\ }\ E_A\ \xrightarrow{\ u_A\ }\ T .
\end{equation}
\end{proposition}

\begin{proof}
The first inclusion is reflexivity: $s\sim_A s$ for every $s$, since $p(s)=p(s)$ and
$r_\alpha(s)=r_\alpha(s)$. The second holds because the defining condition \eqref{eq:simA} contains
the clause $p(s)=p(s')$ as a conjunct; hence $s\sim_A s'$ implies $s\sim_p s'$. The factorization
\eqref{eq:tower} is \eqref{eq:factorization}.
\end{proof}

The reading of \eqref{eq:ordering} must be got right, because the direction of the inclusions
encodes the direction of every universal arrow in the paper. A \emph{smaller} equivalence relation
identifies \emph{fewer} pairs, hence corresponds to a \emph{finer} partition and to \emph{more}
retained information. Thus \eqref{eq:ordering} says: equality on $S$ retains everything;
$\sim_A$ retains at most that, and at least the readout; $\sim_p$ retains the least of the three.
The completion $E_A$ refines the readout partition, but---this is the whole point of minimality---it
need not refine it all the way down to individual source states.

\subsection{What completion adds above a readout value}

\begin{proposition}[Completed fibers]
\label{prop:fibers}
For every $t\in T$ there is a canonical bijection
\begin{equation}
\label{eq:fibers}
u_A^{-1}(t)\ \cong\ R_A(F_t)=\{\,R_A(s)\ :\ s\in F_t=p^{-1}(t)\,\}\ \subseteq\ V_A .
\end{equation}
\end{proposition}

\begin{proof}
Consider the second-coordinate projection $\mathrm{pr}_{V_A}\colon T\times V_A\to V_A$ and restrict
it to $u_A^{-1}(t)$.

\emph{Its image is $R_A(F_t)$.} An element $x\in u_A^{-1}(t)$ lies in $E_A=\Image(j_A)$, so
$x=(p(s),R_A(s))$ for some $s\in S$; and $u_A(x)=t$ forces $p(s)=t$, i.e.\ $s\in F_t$. Hence
$\mathrm{pr}_{V_A}(x)=R_A(s)\in R_A(F_t)$. Conversely, given $v=R_A(s)$ with $s\in F_t$, the point
$x=(t,v)=j_A(s)$ lies in $E_A$ and satisfies $u_A(x)=t$, so $v$ is attained.

\emph{It is injective on $u_A^{-1}(t)$.} If $x,x'\in u_A^{-1}(t)$ have the same $V_A$-coordinate,
then they also share the $T$-coordinate, namely $t$; being elements of a product with all
coordinates equal, $x=x'$.

Thus $\mathrm{pr}_{V_A}$ restricts to a bijection $u_A^{-1}(t)\to R_A(F_t)$, and this bijection was
constructed without any choice.
\end{proof}

\Cref{prop:fibers} answers exactly the question that ``completion'' invites one to ask badly. What
sits above a readout value $t$ in $E_A$ is \emph{not} a preferred source state in $F_t$, and not an
arbitrary enrichment of $t$ by extra labels. It is the set of \emph{realized relation signatures}
over $t$: the distinct values that the nominated family actually takes on the fiber. Two source
states in $F_t$ that agree on every $r_\alpha$ contribute the \emph{same} completed state; nothing in
the construction selects between them, and nothing in the construction can. We record this as a
standing prohibition, because it is the failure mode that a naive completion invites.

\begin{remark}[No selector, no section, no representative]
\label{rem:no-selector}
The construction of $E_A$ uses no selection function on the fibers of $p$, no section of $p$, and no
canonical representative of $F_t$; it requires no choice principle beyond forming the image of a
function, and it produces none of these objects either. This should be stated precisely, since the
point is about what the construction \emph{supplies}, not about what is set-theoretically possible.

If one separately chooses a representative of each $\sim_A$-class---a choice that in general requires
an appeal to a choice principle, and which is available when such a choice can be made---then the
resulting set of chosen representatives is in bijection with $E_A$, by \cref{prop:quotient}. Such a
set is a \emph{noncanonical presentation} of the completion. It carries no invariant information
beyond what $E_A$ already carries: the bijection depends on the choice, different choices give
different presentations of the same quotient, and nothing in the completion distinguishes one of them.
Crucially, no such presentation is supplied or privileged by the universal construction. The object
$E_A$ is defined without any choice whatever, as an image; and terminality
(\cref{thm:terminal}) yields only the surjections $\phi_{\mathcal E}\colon E\to E_A$ \emph{onto} the
completion. It furnishes no inverse to $e_A$, no canonical splitting, and no canonical representative
selector. Whatever a representative-based presentation may be used for, it is extra structure imposed
from outside, not a consequence of the theorem.
\end{remark}

\section{The category of relation-sufficient extensions}
\label{sec:category}

We now fix the data $(S,T,p,\{r_\alpha\}_{\alpha\in A})$ for the remainder of the paper and make
\cref{def:problem} precise as a category.

\begin{definition}[Relation-sufficient extension]
\label{def:sufficient}
A \emph{relation-sufficient extension} of $p$ (relative to the nominated family) is a tuple
\begin{equation}
\label{eq:suff-ext}
\mathcal E=\bigl(E,\ e,\ u,\ (\rho_\alpha^E)_{\alpha\in A}\bigr)
\end{equation}
consisting of a set $E$, a \textbf{surjection} $e\colon S\twoheadrightarrow E$, and maps
$u\colon E\to T$ and $\rho_\alpha^E\colon E\to V_\alpha$ ($\alpha\in A$) satisfying
\begin{equation}
\label{eq:suff-factorization}
p=u\circ e,
\qquad
r_\alpha=\rho_\alpha^E\circ e
\qquad(\alpha\in A).
\end{equation}
\end{definition}

Two features of \cref{def:sufficient} deserve comment.

First, the factorizations \eqref{eq:suff-factorization} are exactly requirements (1) and (2) of
\cref{def:problem}: the readout survives, and every nominated relation is a function of the
completed state.

Second, surjectivity of $e$ is not a technical convenience but a substantive part of the definition,
and it is what makes the universal property come out uniquely. A completed state that no source
state realizes is not a state of the system; it is a spurious element. If unrealized elements were
permitted, then given any sufficient $E$ one could adjoin an arbitrary point $\ast$ with arbitrary
values $u(\ast)$ and $\rho_\alpha^E(\ast)$ (subject to nothing, since \eqref{eq:suff-factorization}
constrains only the image of $e$), and the map to $E_A$ would no longer be determined: the value
$\phi(\ast)$ would be unconstrained by $e_A=\phi\circ e$. Uniqueness in \cref{thm:terminal} would
fail. Surjectivity excludes exactly this, and it does so for the honest reason: a completion consists
of states \emph{realized} by source states.

\begin{definition}[Morphism]
\label{def:morphism}
Let $\mathcal E=(E,e,u,(\rho_\alpha^E))$ and $\mathcal E'=(E',e',u',(\rho_\alpha^{E'}))$ be
relation-sufficient extensions. A \emph{morphism} $f\colon\mathcal E\to\mathcal E'$ is a function
$f\colon E\to E'$ satisfying
\begin{equation}
\label{eq:morphism-over-S}
e'=f\circ e .
\end{equation}
\end{definition}

Thus a morphism is a map \emph{under $S$}: it is compatible with the two completion maps out of the
source. We now check that this single condition already forces compatibility with all the structure
maps.

\begin{lemma}[Structure maps are automatically preserved]
\label{lem:structure-preserved}
If $f\colon\mathcal E\to\mathcal E'$ is a morphism in the sense of \cref{def:morphism}, then
\begin{equation}
\label{eq:structure-preserved}
u'\circ f=u,
\qquad
\rho_\alpha^{E'}\circ f=\rho_\alpha^E
\qquad(\alpha\in A).
\end{equation}
\end{lemma}

\begin{proof}
Let $x\in E$. Since $e$ is surjective, choose $s\in S$ with $x=e(s)$. Then, using
\eqref{eq:morphism-over-S} and the factorizations \eqref{eq:suff-factorization} for $\mathcal E'$
and then for $\mathcal E$,
\[
u'\bigl(f(x)\bigr)=u'\bigl(f(e(s))\bigr)=u'\bigl(e'(s)\bigr)=p(s)=u\bigl(e(s)\bigr)=u(x),
\]
and likewise, for each $\alpha\in A$,
\[
\rho_\alpha^{E'}\bigl(f(x)\bigr)
=\rho_\alpha^{E'}\bigl(e'(s)\bigr)
=r_\alpha(s)
=\rho_\alpha^{E}\bigl(e(s)\bigr)
=\rho_\alpha^{E}(x).
\]
Both computations are independent of the chosen preimage $s$, since their left-hand sides do not
mention $s$. As $x$ was arbitrary, \eqref{eq:structure-preserved} holds.
\end{proof}

We display \eqref{eq:structure-preserved} rather than merely deducing it, because it records what a
morphism of completions actually preserves: the readout and every nominated relation value. What is
\emph{not} imposed is any compatibility beyond that---no metric, no ordering, no choice of labels.

\begin{proposition}[This is a category]
\label{prop:category}
Relation-sufficient extensions of $p$ relative to $\{r_\alpha\}_{\alpha\in A}$, with the morphisms of
\cref{def:morphism}, form a category, which we denote $\mathbf{Suff}(p;\{r_\alpha\})$.
\end{proposition}

\begin{proof}
For $\mathcal E=(E,e,u,(\rho_\alpha^E))$, the identity function $\id_E$ satisfies $e=\id_E\circ e$,
so it is a morphism $\mathcal E\to\mathcal E$. If $f\colon\mathcal E\to\mathcal E'$ and
$g\colon\mathcal E'\to\mathcal E''$ are morphisms, then
$(g\circ f)\circ e=g\circ(f\circ e)=g\circ e'=e''$, so $g\circ f\colon\mathcal E\to\mathcal E''$ is a
morphism. Associativity of composition and the unit laws are inherited from composition of functions.
\end{proof}

We add no further categorical apparatus; none is needed.

\section{The canonical minimal-completion theorem}
\label{sec:main}

\begin{theorem}[Canonical minimal completion]
\label{thm:terminal}
The tuple
\begin{equation}
\label{eq:EA-object}
\mathcal E_A=\bigl(E_A,\ e_A,\ u_A,\ (\rho_\alpha)_{\alpha\in A}\bigr)
\end{equation}
of \cref{def:EA} is an object of $\mathbf{Suff}(p;\{r_\alpha\})$, and it is \textbf{terminal}: for
every relation-sufficient extension $\mathcal E=(E,e,u,(\rho_\alpha^E))$ there exists a
\textbf{unique} morphism
\begin{equation}
\label{eq:phiE-direction}
\phi_{\mathcal E}\colon\ \mathcal E\longrightarrow\mathcal E_A,
\qquad\text{that is,}\qquad
\phi_{\mathcal E}\colon E\longrightarrow E_A .
\end{equation}
It is given by the explicit formula
\begin{equation}
\label{eq:phiE}
\phi_{\mathcal E}(x)=\Bigl(u(x),\ \bigl(\rho_\alpha^E(x)\bigr)_{\alpha\in A}\Bigr),
\qquad x\in E,
\end{equation}
and $\phi_{\mathcal E}$ is surjective.
\end{theorem}

Before the proof, we emphasize the direction, since reversing it inverts the meaning of the theorem.
The universal map runs \textbf{from an arbitrary sufficient extension to $E_A$}, not from $E_A$ to
it. $E_A$ is therefore the \emph{terminal}---the \emph{coarsest}---sufficient extension: every
sufficient extension maps \emph{onto} it, and so retains at least as much information as it does.
$E_A$ is not an initial object, not a universal cover, and not the finest extension; the finest
sufficient extension is $S$ itself (with $e=\id_S$), which is precisely the object one is trying not
to be forced into.

\begin{proof}
\emph{$\mathcal E_A$ is an object.} $e_A$ is surjective and the factorizations
\eqref{eq:factorization} hold, by \cref{prop:factorization}.

Now fix a relation-sufficient extension $\mathcal E=(E,e,u,(\rho_\alpha^E))$ and define
$\phi_{\mathcal E}$ by the formula \eqref{eq:phiE}. We establish the required properties in turn.

\smallskip
\emph{(i) The formula is well defined as a function $E\to T\times V_A$.} For each $x\in E$ the
value $u(x)\in T$ is determined, and for each $\alpha\in A$ the value $\rho_\alpha^E(x)\in V_\alpha$
is determined; the assignment $\alpha\mapsto\rho_\alpha^E(x)$ is therefore an element of the
set-theoretic product $V_A=\prod_{\alpha\in A}V_\alpha$. Hence \eqref{eq:phiE} assigns to each
$x\in E$ exactly one element of $T\times V_A$. No choice is involved, and no representative of $x$
is used: the formula depends only on $x$ through the given structure maps of $\mathcal E$.

\smallskip
\emph{(ii) The image lies in $E_A$, and $e_A=\phi_{\mathcal E}\circ e$.} Let $x\in E$. By
surjectivity of $e$ there is $s\in S$ with $x=e(s)$. Then, using
\eqref{eq:suff-factorization},
\begin{equation}
\label{eq:phiE-computation}
\phi_{\mathcal E}(x)
=\phi_{\mathcal E}\bigl(e(s)\bigr)
=\Bigl(u\bigl(e(s)\bigr),\ \bigl(\rho_\alpha^E(e(s))\bigr)_{\alpha\in A}\Bigr)
=\Bigl(p(s),\ \bigl(r_\alpha(s)\bigr)_{\alpha\in A}\Bigr)
=j_A(s)
=e_A(s).
\end{equation}
The right-hand side lies in $\Image(j_A)=E_A$ by \cref{def:EA}. Since $x\in E$ was arbitrary,
$\phi_{\mathcal E}(E)\subseteq E_A$, so $\phi_{\mathcal E}$ may be regarded as a map $E\to E_A$; and
reading \eqref{eq:phiE-computation} from the ends, $\phi_{\mathcal E}(e(s))=e_A(s)$ for every $s\in
S$, i.e.
\begin{equation}
\label{eq:phiE-over-S}
e_A=\phi_{\mathcal E}\circ e .
\end{equation}
Note that the computation \eqref{eq:phiE-computation} is legitimate for \emph{every} choice of
preimage $s$ of $x$, and its left-hand side does not depend on $s$; so no consistency question
arises.

\smallskip
\emph{(iii) It is a morphism, and preserves all structure maps.} Equation \eqref{eq:phiE-over-S} is
precisely condition \eqref{eq:morphism-over-S} of \cref{def:morphism}, so $\phi_{\mathcal E}$ is a
morphism $\mathcal E\to\mathcal E_A$. Preservation of the structure maps then follows from
\cref{lem:structure-preserved}, but in the present case it is also visible directly from
\eqref{eq:phiE}: applying the coordinate projections $u_A$ and $\rho_\alpha$ to
$\phi_{\mathcal E}(x)=\bigl(u(x),(\rho_\beta^E(x))_\beta\bigr)$ gives
\begin{equation}
\label{eq:phiE-structure}
u_A\circ\phi_{\mathcal E}=u,
\qquad
\rho_\alpha\circ\phi_{\mathcal E}=\rho_\alpha^E
\qquad(\alpha\in A).
\end{equation}

\smallskip
\emph{(iv) Surjectivity.} Let $y\in E_A$. By \cref{prop:factorization} the map $e_A$ is surjective,
so $y=e_A(s)$ for some $s\in S$. By \eqref{eq:phiE-over-S}, $y=\phi_{\mathcal E}(e(s))$, and
$e(s)\in E$. Hence $\phi_{\mathcal E}$ is onto $E_A$.

\smallskip
\emph{(v) Uniqueness.} Suppose $\psi\colon\mathcal E\to\mathcal E_A$ is any morphism, i.e.\ any
function $\psi\colon E\to E_A$ with $e_A=\psi\circ e$. Let $x\in E$ and choose $s\in S$ with
$x=e(s)$. Then
\[
\psi(x)=\psi\bigl(e(s)\bigr)=e_A(s)=\phi_{\mathcal E}\bigl(e(s)\bigr)=\phi_{\mathcal E}(x),
\]
using $e_A=\psi\circ e$ for the second equality and \eqref{eq:phiE-over-S} for the third. Since $x$
was arbitrary, $\psi=\phi_{\mathcal E}$. (This is exactly where surjectivity of $e$ is consumed: it
guarantees that the condition $e_A=\psi\circ e$ pins down $\psi$ at \emph{every} point of $E$, not
merely on a subset.)
\end{proof}

\begin{corollary}[Uniqueness up to unique isomorphism]
\label{cor:unique}
Any two terminal objects of the category $\mathbf{Suff}(p;\{r_\alpha\})$ are isomorphic
by a unique isomorphism, and this isomorphism commutes with the completion maps from $S$
and with all of the structure maps. In particular, any relation-sufficient extension that
satisfies the universal property of \cref{thm:terminal} is canonically isomorphic to
$\mathcal E_A$.
\end{corollary}

\begin{proof}
Let $\mathcal E,\mathcal E'$ be terminal. Terminality of $\mathcal E'$ gives a unique morphism
$f\colon\mathcal E\to\mathcal E'$; terminality of $\mathcal E$ gives a unique morphism
$g\colon\mathcal E'\to\mathcal E$. Then $g\circ f\colon\mathcal E\to\mathcal E$ is a morphism, as is
$\id_{\mathcal E}$; by terminality of $\mathcal E$ there is only one morphism
$\mathcal E\to\mathcal E$, so $g\circ f=\id_E$. Symmetrically $f\circ g=\id_{E'}$. Hence $f$ is an
isomorphism, and it is unique because morphisms into a terminal object are unique. Compatibility with
the completion maps is \eqref{eq:morphism-over-S}; compatibility with the structure maps is
\cref{lem:structure-preserved}.
\end{proof}

\begin{corollary}[Coarsest sufficiency, elementary form]
\label{cor:coarsest}
Let $\mathcal E=(E,e,u,(\rho_\alpha^E))$ be any relation-sufficient extension, and define its induced
equivalence relation by
\begin{equation}
\label{eq:sim-e}
s\sim_e s'\quad:\Longleftrightarrow\quad e(s)=e(s').
\end{equation}
Then
\begin{equation}
\label{eq:coarsest}
\sim_e\ \subseteq\ \sim_A .
\end{equation}
Equivalently: every sufficient extension separates at least every pair that $E_A$ separates, while
$E_A$ separates no pair that sufficiency does not force it to separate.
\end{corollary}

\begin{proof}
Suppose $e(s)=e(s')$. Applying $u$ and using $p=u\circ e$,
\[
p(s)=u\bigl(e(s)\bigr)=u\bigl(e(s')\bigr)=p(s') .
\]
Applying $\rho_\alpha^E$ and using $r_\alpha=\rho_\alpha^E\circ e$, for each $\alpha\in A$,
\[
r_\alpha(s)=\rho_\alpha^E\bigl(e(s)\bigr)=\rho_\alpha^E\bigl(e(s')\bigr)=r_\alpha(s') .
\]
By \eqref{eq:simA} this is precisely $s\sim_A s'$.
\end{proof}

\Cref{cor:coarsest} and \cref{thm:terminal} are two faces of one statement. Terminality is the
categorical form; the inclusion \eqref{eq:coarsest} is the elementary, checkable form. The
inclusion also settles the direction question once and for all: since $\sim_e$ is contained in
$\sim_A$, the partition induced by $e$ \emph{refines} the partition induced by $e_A$, so there is a
well-defined map from the finer quotient onto the coarser one---from $E$ onto $E_A$, and not back.

\begin{remark}[The sense in which the completion is a construction]
\label{rem:construction-scope}
\Cref{thm:terminal} constructs the coarsest relation-sufficient extension in the paper's category
$\mathbf{Suff}(p;\{r_\alpha\})$ of set-theoretic source quotients: the object $E_A$ exists, is
terminal, and is unique up to unique isomorphism, and those claims are unconditional. What the
theorem builds it builds from the nominated relations themselves. The joint map
$j_A=(p,R_A)$ takes the $r_\alpha$ as given input, so the construction presupposes the relation
data rather than producing it; in particular it furnishes \emph{no} procedure for recovering the
$r_\alpha$ from the readout $p$ when they fail to descend---if such a procedure existed the
relations would already be functions of $p$, and the completion would be $T$ itself. Nor does
terminality establish that $E_A$ is experimentally accessible, measurable, smooth, geometric, or
realized as a physical state space independent of its defining coordinates; $E_A$ is, as
constructed, a quotient of $S$ recording exactly the nominated values. A concrete realization must
independently supply and justify the relation data, and any additional geometric or physical
structure carried by $E_A$, on its own terms. That realization question is not settled here; it is
deferred to the companion manuscript on geometric realization discussed in \cref{sec:discussion}.
\end{remark}

\begin{remark}[What ``minimal'' means, and what it does not]
\label{rem:minimality}
Minimality here means: \emph{least retained distinction among sufficient extensions}. It is a
statement about the ordering of equivalence relations by inclusion, equivalently about the ordering
of partitions by refinement. It does \emph{not} mean minimum cardinality selected without regard to
structure---the one-point set is smaller than $E_A$ and is not, in general, sufficient---and it does
\emph{not} mean that $E_A$ is a small or economical description of $S$ in any other sense. Above all,
minimal completion is \emph{not} reconstruction of $S$: by \cref{prop:ordering} the map $e_A$ need
not be injective, and \cref{sec:reconstruction} determines exactly when it is.
\end{remark}

\section{Redundant and indispensable relation coordinates}
\label{sec:redundancy}

The nominated family may be wasteful: one of its members may already be determined by the readout
together with the others, in which case its presence in the list changes nothing. We characterize
this exactly.

Throughout this section fix $\alpha\in A$ and write
\begin{equation}
\label{eq:A-minus}
A_{-\alpha}:=A\setminus\{\alpha\},
\qquad
R_{A_{-\alpha}}=(r_\beta)_{\beta\in A_{-\alpha}}\colon S\to V_{A_{-\alpha}} .
\end{equation}

\begin{definition}[Redundancy and indispensability]
\label{def:redundant}
The relation $r_\alpha$ is \emph{redundant relative to $p$ and the remaining relations} if it factors
through their joint map, i.e.\ if there exists a function
$f_\alpha\colon T\times V_{A_{-\alpha}}\to V_\alpha$ with
\begin{equation}
\label{eq:redundant}
r_\alpha=f_\alpha\circ\bigl(p,\ R_{A_{-\alpha}}\bigr).
\end{equation}
It is \emph{indispensable} if no such factorization exists.
\end{definition}

\begin{theorem}[Indispensability criterion]
\label{thm:indispensable}
The following are equivalent.
\begin{enumerate}
\item\label{it:indisp} $r_\alpha$ is indispensable.
\item\label{it:witness} There exist $s,s'\in S$ with
\begin{equation}
\label{eq:witness}
p(s)=p(s'),
\qquad
r_\beta(s)=r_\beta(s')\quad\text{for every }\beta\in A_{-\alpha},
\qquad
r_\alpha(s)\neq r_\alpha(s').
\end{equation}
\end{enumerate}
\end{theorem}

\begin{proof}
We prove the contrapositive equivalence: $r_\alpha$ is redundant if and only if no pair
\eqref{eq:witness} exists, i.e.\ if and only if $r_\alpha$ is constant on every fiber of the joint
map $g:=(p,R_{A_{-\alpha}})\colon S\to T\times V_{A_{-\alpha}}$. This is the descent statement of
Paper~1 \cite{Paper1} applied to $g$; we give the argument in full.

\emph{Redundant $\Rightarrow$ constant on fibers of $g$.} Suppose \eqref{eq:redundant} holds and
suppose $g(s)=g(s')$, i.e.\ $p(s)=p(s')$ and $r_\beta(s)=r_\beta(s')$ for all $\beta\neq\alpha$.
Applying $f_\alpha$ to the equal values $g(s)=g(s')$ gives
$r_\alpha(s)=f_\alpha(g(s))=f_\alpha(g(s'))=r_\alpha(s')$. Hence no pair as in \eqref{eq:witness}
can exist.

\emph{Constant on fibers of $g$ $\Rightarrow$ redundant.} Suppose that $g(s)=g(s')$ implies
$r_\alpha(s)=r_\alpha(s')$. Define $f_\alpha$ on the image of $g$ by
\begin{equation}
\label{eq:falpha-def}
f_\alpha\bigl(g(s)\bigr):=r_\alpha(s)\qquad(s\in S),
\end{equation}
The prescription \eqref{eq:falpha-def} is well defined on $\Image(g)$ precisely by the assumed
fiber-constancy: if $g(s)=g(s')$ then the two prescribed values $r_\alpha(s)$ and $r_\alpha(s')$
agree. Every point of $\Image(g)$ is of the form $g(s)$, so $f_\alpha$ is defined on all of
$\Image(g)$. It remains to extend $f_\alpha$ to the whole of $T\times V_{A_{-\alpha}}$, and we treat
the two cases explicitly.

\emph{Case $S\neq\varnothing$.} Fix any $s_0\in S$ and define $f_\alpha(y):=r_\alpha(s_0)$ for every
$y\in\bigl(T\times V_{A_{-\alpha}}\bigr)\setminus\Image(g)$. This is legitimate because
$r_\alpha(s_0)$ is an element of $V_\alpha$, so $V_\alpha\neq\varnothing$ and a value is available.
The extension changes nothing on $\Image(g)$, where \eqref{eq:falpha-def} already fixes $f_\alpha$.

\emph{Case $S=\varnothing$.} Then $T=\varnothing$: the readout $p\colon S\to T$ is surjective by
standing hypothesis, and the image of the empty map is empty, so $T=p(S)=\varnothing$. Hence the
product $T\times V_{A_{-\alpha}}=\varnothing$ as well, and the unique map
$\varnothing\to V_\alpha$---the empty map---serves as $f_\alpha$. It satisfies
$r_\alpha=f_\alpha\circ g$ vacuously as an identity of maps out of $S=\varnothing$, both sides being
the empty map. So $r_\alpha$ is redundant in this case, which is consistent: no witness pair
\eqref{eq:witness} can exist in an empty source, so $r_\alpha$ is not indispensable, exactly as the
equivalence requires. (The two sides of the criterion agree here; the case is not merely set aside.)

In either case $f_\alpha$ is defined on all of $T\times V_{A_{-\alpha}}$ and satisfies
$r_\alpha=f_\alpha\circ g$, which is \eqref{eq:redundant}.

Finally, ``no pair as in \eqref{eq:witness}'' is by inspection the same statement as ``$g(s)=g(s')$
implies $r_\alpha(s)=r_\alpha(s')$'', since the first two clauses of \eqref{eq:witness} say
$g(s)=g(s')$ and the third says $r_\alpha(s)\neq r_\alpha(s')$. Negating both sides of the
equivalence just proved yields \labelcref{it:indisp}$\Leftrightarrow$\labelcref{it:witness}.
\end{proof}

The criterion is a genuine separation test, and it is the reason a long list of nominated relations
is not thereby an economical one. A coordinate earns its place only by separating some pair of source
states that the readout and all the \emph{other} coordinates identify.

\begin{corollary}[Strict coarsening on deletion]
\label{cor:deletion}
Let $\sim_A$ and $\sim_{A_{-\alpha}}$ be the equivalence relations \eqref{eq:simA} induced by the
full family and by the family with $r_\alpha$ deleted, and let $\pi\colon E_A\to E_{A_{-\alpha}}$ be
the canonical surjection of \cref{prop:refinement} below (with $B=A$). Then always
$\sim_A\ \subseteq\ \sim_{A_{-\alpha}}$, and the following are equivalent:
\begin{enumerate}
\item the inclusion is \emph{strict}: $\sim_A\ \subsetneq\ \sim_{A_{-\alpha}}$;
\item $\pi$ is not injective, i.e.\ deleting $r_\alpha$ \emph{strictly coarsens} the completion;
\item $r_\alpha$ is indispensable.
\end{enumerate}
\end{corollary}

\begin{proof}
The inclusion $\sim_A\subseteq\sim_{A_{-\alpha}}$ is immediate from \eqref{eq:simA}, since the
defining condition for $\sim_A$ is the condition for $\sim_{A_{-\alpha}}$ conjoined with the extra
clause $r_\alpha(s)=r_\alpha(s')$.

\emph{(1) $\Leftrightarrow$ (3).} Strictness means: there exist $s,s'$ with $s\sim_{A_{-\alpha}}s'$
but not $s\sim_A s'$. Unwinding \eqref{eq:simA}, this says $p(s)=p(s')$ and $r_\beta(s)=r_\beta(s')$
for all $\beta\neq\alpha$, while $r_\alpha(s)\neq r_\alpha(s')$---which is exactly the witness
\eqref{eq:witness}. By \cref{thm:indispensable} this holds iff $r_\alpha$ is indispensable.

\emph{(1) $\Leftrightarrow$ (2).} By \cref{prop:quotient}, $E_A\cong S/\!\sim_A$ and
$E_{A_{-\alpha}}\cong S/\!\sim_{A_{-\alpha}}$, and under these identifications $\pi$ is the canonical
map $S/\!\sim_A\ \to\ S/\!\sim_{A_{-\alpha}}$ sending $[s]_A\mapsto[s]_{A_{-\alpha}}$, which is
well defined by the inclusion just proved. This map is injective iff
$[s]_{A_{-\alpha}}=[s']_{A_{-\alpha}}$ implies $[s]_A=[s']_A$, i.e.\ iff
$\sim_{A_{-\alpha}}\ \subseteq\ \sim_A$, i.e.\ (given the reverse inclusion) iff the two relations are
equal. Its failure is precisely strictness.
\end{proof}

\begin{remark}[Irredundancy is not uniqueness of coordinates]
\label{rem:not-unique-coords}
A family in which every member is indispensable is \emph{irredundant relative to the nomination}.
This is a statement about the given list, and it confers no uniqueness on the list. Different
irredundant families may induce the very same equivalence relation $\sim_A$ and hence, by
\cref{prop:quotient}, the very same completion---\cref{prop:reparam} below exhibits an entire class
of such rewritings. An irredundant list of relation coordinates must therefore not be confused with a
preferred coordinate system on the completion.

The correct statement of the coordinate situation is a relative one, and we make it precisely.
\emph{Relative to the fixed nominated family} $\{r_\alpha\}_{\alpha\in A}$, the completion $E_A$ does
carry canonical structure maps: the readout projection $u_A\colon E_A\to T$ and the relation
coordinates $\rho_\alpha\colon E_A\to V_\alpha$ of \eqref{eq:structure-maps}, which are part of the
data of the terminal object and are determined by it. What $E_A$ does \emph{not} possess is a
family-independent, intrinsically unique coordinate presentation: a system of coordinates singled out
by the quotient alone, without reference to a nomination. Faithful reparameterizations
(\cref{prop:reparam}) and distinct irredundant families may present one and the same quotient in
different coordinates, and the quotient does not choose among them. The invariant is $E_A$ together
with its structure maps for the given nomination; the coordinates are a presentation relative to that
nomination.
\end{remark}

\section{Nested families, refinement maps, and reparameterization}
\label{sec:refinement}

\subsection{Nested families}

Let $A\subseteq B$ be sets of indices, the $A$-relations being a subfamily of the $B$-relations
(i.e.\ the relations indexed by $A$ are the same functions in both families). Then, directly from
\eqref{eq:simA},
\begin{equation}
\label{eq:nested-order}
\sim_B\ \subseteq\ \sim_A ,
\end{equation}
since agreement of \emph{all} $B$-relations (and the readout) entails agreement of the sub-collection
indexed by $A$ (and the readout). In words: a larger nomination is at least as demanding, hence
yields an at least as fine completion.

\begin{proposition}[Canonical refinement maps]
\label{prop:refinement}
Let $A\subseteq B$. There is a unique map
\begin{equation}
\label{eq:piBA}
\pi_{BA}\colon E_B\longrightarrow E_A
\qquad\text{with}\qquad
e_A=\pi_{BA}\circ e_B ,
\end{equation}
and it is surjective. It satisfies $u_A\circ\pi_{BA}=u_B$ and
$\rho_\alpha^{(A)}\circ\pi_{BA}=\rho_\alpha^{(B)}$ for every $\alpha\in A$, where the superscripts
indicate the completion in which the structure map lives. Moreover $\pi_{AA}=\id_{E_A}$, and for
$A\subseteq B\subseteq C$,
\begin{equation}
\label{eq:cocycle}
\pi_{CA}=\pi_{BA}\circ\pi_{CB}.
\end{equation}
\end{proposition}

\begin{proof}
\emph{Existence and uniqueness.} We exhibit $E_B$ as a relation-sufficient extension \emph{for the
$A$-family} and invoke \cref{thm:terminal}. Indeed, $e_B\colon S\twoheadrightarrow E_B$ is
surjective; $u_B\colon E_B\to T$ satisfies $p=u_B\circ e_B$; and for each $\alpha\in A\subseteq B$
the structure map $\rho_\alpha^{(B)}\colon E_B\to V_\alpha$ satisfies
$r_\alpha=\rho_\alpha^{(B)}\circ e_B$. Hence
$\bigl(E_B,e_B,u_B,(\rho_\alpha^{(B)})_{\alpha\in A}\bigr)$ is an object of
$\mathbf{Suff}(p;\{r_\alpha\}_{\alpha\in A})$. \Cref{thm:terminal} supplies a unique morphism from it
to the terminal object $\mathcal E_A$, i.e.\ a unique map $\pi_{BA}\colon E_B\to E_A$ with
$e_A=\pi_{BA}\circ e_B$, and asserts that it is surjective. Compatibility with $u$ and with the
$\rho_\alpha$, $\alpha\in A$, is \cref{lem:structure-preserved} (equivalently
\eqref{eq:phiE-structure}).

\emph{Identity.} For $B=A$, the identity $\id_{E_A}$ satisfies $e_A=\id_{E_A}\circ e_A$; by the
uniqueness just established, $\pi_{AA}=\id_{E_A}$.

\emph{Cocycle law.} Let $A\subseteq B\subseteq C$. Both $\pi_{CA}$ and $\pi_{BA}\circ\pi_{CB}$ are
maps $E_C\to E_A$. Compose each with the surjection $e_C$:
\[
\pi_{CA}\circ e_C=e_A,
\qquad
(\pi_{BA}\circ\pi_{CB})\circ e_C
=\pi_{BA}\circ(\pi_{CB}\circ e_C)
=\pi_{BA}\circ e_B
=e_A .
\]
Two maps out of $E_C$ that agree after precomposition with the surjection $e_C$ agree on all of
$E_C$: given $x\in E_C$, write $x=e_C(s)$ and evaluate. Hence \eqref{eq:cocycle}. (Alternatively,
both maps are morphisms $\mathcal E_C\to\mathcal E_A$ in $\mathbf{Suff}(p;\{r_\alpha\}_{\alpha\in
A})$, and \cref{thm:terminal} says there is only one.)
\end{proof}

Nested nominations therefore generate a coherent tower of completions,
\begin{equation}
\label{eq:tower-nested}
S\ \longrightarrow\ E_C\ \xrightarrow{\ \pi_{CB}\ }\ E_B\ \xrightarrow{\ \pi_{BA}\ }\ E_A\
\xrightarrow{\ u_A\ }\ T
\qquad(A\subseteq B\subseteq C),
\end{equation}
in which each arrow forgets exactly the distinctions that the smaller nomination does not require.
This is a hierarchy of \emph{information content}. Its arrows run from finer to coarser, and by
\cref{thm:terminal} each is uniquely determined by the pair of nominations.

\begin{remark}[Strictness is not automatic]
\label{rem:not-strict}
The completion for a larger family is at least as fine as that for a smaller one, and \emph{may} be
strictly finer, but need not be. If every relation indexed by $B\setminus A$ is redundant relative to
$p$ and the $A$-relations, then $\sim_B\ =\ \sim_A$ by the argument of \cref{cor:deletion}, and
$\pi_{BA}$ is a bijection. Adding relations to the nomination adds information to the completion
precisely to the extent that the added relations separate pairs that the existing nomination
identifies.
\end{remark}

\subsection{Faithful reparameterization}

The completion depends on the \emph{distinctions} carried by the nominated relations, not on the
labels in which their values happen to be recorded. We make this precise.

\begin{proposition}[Invariance under faithful reparameterization]
\label{prop:reparam}
For each $\alpha\in A$ let $f_\alpha\colon V_\alpha\to W_\alpha$ be a map whose restriction to the
realized subset $r_\alpha(S)\subseteq V_\alpha$ is injective, and set
\begin{equation}
\label{eq:reparam}
r'_\alpha:=f_\alpha\circ r_\alpha\colon S\to W_\alpha .
\end{equation}
Let $\sim_A$ and $\sim'_A$ be the equivalence relations induced by $\{r_\alpha\}$ and by
$\{r'_\alpha\}$ (with the same readout $p$), and let $E_A$, $E'_A$ be the corresponding completions.
Then
\begin{equation}
\label{eq:same-relation}
\sim_A\ =\ \sim'_A ,
\end{equation}
and consequently there is a canonical bijection $E_A\cong E'_A$ commuting with the completion maps
from $S$ and with the projections to $T$:
\begin{equation}
\label{eq:iso-over-T}
u'_A\circ\bigl(E_A\xrightarrow{\ \cong\ }E'_A\bigr)=u_A .
\end{equation}
\end{proposition}

\begin{proof}
\emph{$\sim_A\subseteq\sim'_A$.} Suppose $s\sim_A s'$. Then $p(s)=p(s')$, and for each $\alpha$,
$r_\alpha(s)=r_\alpha(s')$; applying the function $f_\alpha$ to equal arguments gives
$r'_\alpha(s)=f_\alpha(r_\alpha(s))=f_\alpha(r_\alpha(s'))=r'_\alpha(s')$. Hence $s\sim'_A s'$. (Only
well-definedness of $f_\alpha$ is used here; injectivity is not.)

\emph{$\sim'_A\subseteq\sim_A$.} Suppose $s\sim'_A s'$. Then $p(s)=p(s')$, and for each $\alpha$,
$f_\alpha(r_\alpha(s))=f_\alpha(r_\alpha(s'))$. Both $r_\alpha(s)$ and $r_\alpha(s')$ lie in the
realized subset $r_\alpha(S)$, on which $f_\alpha$ is injective by hypothesis; therefore
$r_\alpha(s)=r_\alpha(s')$. As this holds for every $\alpha$, $s\sim_A s'$.

This proves \eqref{eq:same-relation}. By \cref{prop:quotient} applied to each family,
$E_A\cong S/\!\sim_A\ =\ S/\!\sim'_A\ \cong E'_A$, and the composite bijection sends
$e_A(s)\mapsto e'_A(s)$, so it commutes with the completion maps. Since $u_A(e_A(s))=p(s)=u'_A(e'_A(s))$
for every $s$, and $e_A$ is surjective, \eqref{eq:iso-over-T} follows.
\end{proof}

\begin{remark}[The precise form of the invariance]
\label{rem:reparam-precise}
The value spaces $V_\alpha$ and $W_\alpha$ are in general different sets, so $E_A\subseteq T\times
V_A$ and $E'_A\subseteq T\times W_A$ are subsets of different products. \Cref{prop:reparam} therefore
asserts an isomorphism of the \emph{induced quotients over $T$}---equivalently, of the two objects as
completions of the readout---and not an identity, nor an isomorphism internal to a single fixed
category of relation signatures valued in one fixed product. The relation-signature coordinates are a
presentation; the quotient is the invariant. Faithful relabelling of the presentation does not move
the invariant.
\end{remark}

\section{Completion determined by a family of target answers}
\label{sec:targets}

Paper~2 \cite{Paper2} works with a single target $Q$ and its correct-answer map $a\colon S\to A_Q$,
and constructs the coarsest extension of the readout that determines that answer; we will denote that
single-target object by $E_Q$, following Paper~2. We now treat a \emph{family} of targets, and we
choose notation that cannot collide with the single-target object.

Let $\mathcal Q$ be a set of targets, each equipped with a correct-answer map
\begin{equation}
\label{eq:answers}
a_Q\colon S\longrightarrow A_Q,\qquad Q\in\mathcal Q ,
\end{equation}
and define
\begin{equation}
\label{eq:EQfam}
j_{\mathcal Q}\colon S\to T\times\prod_{Q\in\mathcal Q}A_Q,
\qquad
j_{\mathcal Q}(s)=\Bigl(p(s),\ \bigl(a_Q(s)\bigr)_{Q\in\mathcal Q}\Bigr),
\qquad
E_{\mathcal Q}:=\Image\bigl(j_{\mathcal Q}\bigr),
\end{equation}
with induced surjection $e_{\mathcal Q}\colon S\twoheadrightarrow E_{\mathcal Q}$, projection
$u_{\mathcal Q}\colon E_{\mathcal Q}\to T$, and structure maps $\sigma_Q\colon E_{\mathcal Q}\to A_Q$.
Write $\sim_{\mathcal Q}$ for the induced equivalence relation: $s\sim_{\mathcal Q}s'$ iff $p(s)=p(s')$
and $a_Q(s)=a_Q(s')$ for every $Q\in\mathcal Q$.

\begin{proposition}[Target-family completion]
\label{prop:target-completion}
The data $\bigl(E_{\mathcal Q},e_{\mathcal Q},u_{\mathcal Q},(\sigma_Q)_{Q\in\mathcal Q}\bigr)$ is the
terminal object of the category $\mathbf{Suff}\bigl(p;\{a_Q\}_{Q\in\mathcal Q}\bigr)$ of surjective
extensions of the readout on which every answer map $a_Q$ becomes well defined. In particular
$E_{\mathcal Q}\cong S/\!\sim_{\mathcal Q}$, and $E_{\mathcal Q}$ is the coarsest extension of the
readout determining every answer in the family.
\end{proposition}

\begin{proof}
The answer maps $a_Q\colon S\to A_Q$ form a set-indexed family of functions on $S$, exactly as the
$r_\alpha$ do. \Cref{def:EA,def:sufficient,def:morphism} and
\cref{prop:factorization,prop:quotient,thm:terminal,cor:coarsest} were stated and proved for an
arbitrary such family and use no property of the $r_\alpha$ beyond their being functions on $S$.
Substituting $\mathcal Q$ for $A$, $a_Q$ for $r_\alpha$, and $A_Q$ for $V_\alpha$ throughout yields
the assertions verbatim.
\end{proof}

\subsection{Mediation and the canonical surjection}

Now suppose the nominated relation family is \emph{sufficient for the targets}, in the precise sense
that each answer is computable from the readout together with the relation values.

\begin{definition}[Mediation]
\label{def:mediation}
The family $\{r_\alpha\}_{\alpha\in A}$ \emph{mediates} $\mathcal Q$ if for every $Q\in\mathcal Q$
there is a map
\begin{equation}
\label{eq:Phi}
\Phi_Q\colon T\times V_A\longrightarrow A_Q
\qquad\text{with}\qquad
a_Q(s)=\Phi_Q\bigl(p(s),\ R_A(s)\bigr)\quad\text{for all }s\in S .
\end{equation}
\end{definition}

\begin{theorem}[Relation completion versus target completion]
\label{thm:chi}
Suppose $\{r_\alpha\}_{\alpha\in A}$ mediates $\mathcal Q$. Then there is a \emph{unique} map
\begin{equation}
\label{eq:chi}
\chi\colon E_A\longrightarrow E_{\mathcal Q}
\qquad\text{with}\qquad
e_{\mathcal Q}=\chi\circ e_A ,
\end{equation}
and $\chi$ is surjective; explicitly,
$\chi(x)=\bigl(u_A(x),(\Phi_Q(x))_{Q\in\mathcal Q}\bigr)$ for $x\in E_A\subseteq T\times V_A$.
Moreover, $\chi$ is a bijection if and only if
\begin{equation}
\label{eq:chi-iso}
\sim_{\mathcal Q}\ =\ \sim_A ,
\end{equation}
that is, if and only if equality of the readout together with all target answers induces exactly the
same equivalence relation on $S$ as equality of the readout together with all selected relation
values. In general one has only
\begin{equation}
\label{eq:chi-order}
\sim_A\ \subseteq\ \sim_{\mathcal Q} ,
\end{equation}
so the target completion may be \textbf{strictly coarser} than the relation completion.
\end{theorem}

\begin{proof}
\emph{$E_A$ is sufficient for $\mathcal Q$.} Define $\sigma_Q^{A}:=\Phi_Q|_{E_A}\colon E_A\to A_Q$
(the restriction of $\Phi_Q$ to $E_A\subseteq T\times V_A$). For $s\in S$,
\[
\sigma_Q^{A}\bigl(e_A(s)\bigr)=\Phi_Q\bigl(p(s),R_A(s)\bigr)=a_Q(s)
\]
by \eqref{eq:Phi}. Together with the identity $p=u_A\circ e_A$ and the surjectivity of
the map $e_A$, this exhibits the tuple
$\bigl(E_A,e_A,u_A,(\sigma_Q^A)_{Q\in\mathcal Q}\bigr)$ as an object of the category
$\mathbf{Suff}(p;\{a_Q\})$.

\emph{Existence, uniqueness, surjectivity of $\chi$.} By \cref{prop:target-completion},
$E_{\mathcal Q}$ is terminal in that category, so \cref{thm:terminal} supplies a unique morphism
$\chi\colon E_A\to E_{\mathcal Q}$, necessarily satisfying $e_{\mathcal Q}=\chi\circ e_A$ and
necessarily surjective; and by the explicit formula \eqref{eq:phiE} it is given on $x\in E_A$ by
$\chi(x)=\bigl(u_A(x),(\sigma_Q^A(x))_{Q\in\mathcal Q}\bigr)
=\bigl(u_A(x),(\Phi_Q(x))_{Q\in\mathcal Q}\bigr)$.

\emph{The ordering \eqref{eq:chi-order}.} Suppose $s\sim_A s'$. Then $p(s)=p(s')$ and $R_A(s)=R_A(s')$,
so by \eqref{eq:Phi}, $a_Q(s)=\Phi_Q(p(s),R_A(s))=\Phi_Q(p(s'),R_A(s'))=a_Q(s')$ for every $Q$. Hence
$s\sim_{\mathcal Q}s'$. (This is \cref{cor:coarsest} applied to the sufficient extension $E_A$ of the
$\mathcal Q$-problem.)

\emph{Injectivity criterion.} By \cref{prop:quotient} applied to both families, $\chi$ is identified
with the canonical map $S/\!\sim_A\ \to\ S/\!\sim_{\mathcal Q}$, $[s]_A\mapsto[s]_{\mathcal Q}$, which
is well defined by \eqref{eq:chi-order}. Such a map is injective iff
$\sim_{\mathcal Q}\ \subseteq\ \sim_A$; combined with \eqref{eq:chi-order} this is exactly
\eqref{eq:chi-iso}. Since $\chi$ is always surjective, it is a bijection precisely under
\eqref{eq:chi-iso}.
\end{proof}

\begin{remark}[Two different things are being retained]
\label{rem:rel-vs-target}
The distinction encoded in \cref{thm:chi} is substantive and is easy to lose. The \emph{relation}
completion $E_A$ retains the nominated relation \emph{values} themselves. The \emph{target}
completion $E_{\mathcal Q}$ retains only those distinctions among relation values that \emph{affect
the nominated answers}. If two source states carry different relation signatures but the mediating
maps $\Phi_Q$ send both signatures to the same answers for every $Q\in\mathcal Q$, then $E_A$
separates them and $E_{\mathcal Q}$ does not. Hence the strict-coarsening possibility in
\eqref{eq:chi-order} is not an artifact; \cref{sec:u1} exhibits it concretely, where the relation
completion retains an orientation that no answer in the nominated spectral family can detect.
\end{remark}

\section{The gauge-reduction boundary}
\label{sec:gauge}

We return to the raw configuration space $\Omega$ and the physical quotient
\eqref{eq:gauge-quotient}, $q\colon\Omega\to S=\Omega/\!\simeq_{\mathrm{phys}}$. Recall
\cref{def:eligible}: only quantities invariant under $\simeq_{\mathrm{phys}}$ descend to $S$, and
only such quantities are eligible for nomination.

\begin{corollary}[Completion does not restore gauge coordinates]
\label{cor:gauge}
For every nominated family $\{r_\alpha\}_{\alpha\in A}$ of eligible relations, the composite
$e_A\circ q\colon\Omega\to E_A$ is constant on raw physical-equivalence classes:
\begin{equation}
\label{eq:gauge-constancy}
\omega\simeq_{\mathrm{phys}}\omega'
\quad\Longrightarrow\quad
e_A\bigl(q(\omega)\bigr)=e_A\bigl(q(\omega')\bigr).
\end{equation}
\end{corollary}

\begin{proof}
If $\omega\simeq_{\mathrm{phys}}\omega'$ then $q(\omega)=q(\omega')$ by definition of the quotient
map, and applying the function $e_A$ to a single element of $S$ yields a single value.
\end{proof}

Mathematically this is a triviality, and we present it as such: it is a corollary of the order of
operations, not a theorem with content of its own. Its \emph{physical} significance is nevertheless
worth stating precisely, because completion is easily misread as a licence to reintroduce structure
that was deliberately discarded. The precise content of \eqref{eq:gauge-constancy} is that the
composite $e_A\circ q$ is constant on physical-equivalence classes. Consequently, no completion,
however large the nomination, can
\begin{enumerate}
\item distinguish two gauge-equivalent raw configurations $\omega\simeq_{\mathrm{phys}}\omega'$;
\item recover which raw representative $\omega$ was originally supplied, given only the completed
state;
\item restore discarded gauge or coordinate data as \emph{physical} information---that is, as an
invariant of the completed state; or
\item canonically generate a gauge section $\Omega\leftarrow S$ from the completion theorem.
\end{enumerate}
We are careful not to overstate this. We do \emph{not} claim that a representative of a
physical-equivalence class can never be selected in any sense whatever: a separately chosen
gauge-fixing convention, or a section of $q$, may well assign a representative when one exists, and
such devices are routine in practice. The point is about provenance and status. A conventional
gauge-fixing section is \emph{extra, noncanonical structure supplied from outside}. It is not
produced by the completion, is not privileged by terminality, and---this is the substantive
claim---is not a recovery of lost raw information: the information it appears to supply was chosen by
the convention, not extracted from the physics. Two different conventions give two different
representatives of the same physical state, and nothing in $E_A$ adjudicates between them.

Completion therefore refines the readout partition, but only \emph{within} the reduced source $S$;
it cannot resurrect a distinction that the reduction $q$ was defined to destroy. Reduction precedes
completion, and the universal completion does not undo physical equivalence. If one wants such a
distinction back as physics, one must argue that the reduction was wrong---i.e.\ that the quantity
was physical after all, hence eligible in the sense of \cref{def:eligible}---and that is an argument
about physics, not a construction available in this paper.

\section{The topological universal property}
\label{sec:topology}

Assume now that $S$, $T$, and each $V_\alpha$ are topological spaces, and that $p$ and every
$r_\alpha$ are continuous. Give $V_A=\prod_{\alpha\in A}V_\alpha$ the product topology and
$T\times V_A$ the product topology; then $j_A=(p,R_A)$ is continuous, by the universal property of
the product.

\begin{definition}[Topological completion]
\label{def:top-EA}
Equip $E_A$ with the \emph{quotient topology} induced by the surjection
$e_A\colon S\twoheadrightarrow E_A$: a subset $U\subseteq E_A$ is open iff $e_A^{-1}(U)$ is open in
$S$. A \emph{topological relation-sufficient extension} is a relation-sufficient extension
$\mathcal E=(E,e,u,(\rho_\alpha^E))$ in which $E$ is a topological space, $e$ is a \emph{quotient
map}, and $u$ and every $\rho_\alpha^E$ are continuous. Morphisms are continuous maps satisfying
\eqref{eq:morphism-over-S}.
\end{definition}

\begin{proposition}[Continuity of the structure maps]
\label{prop:top-structure}
With the quotient topology of \cref{def:top-EA}, the maps $u_A\colon E_A\to T$ and
$\rho_\alpha\colon E_A\to V_\alpha$ are continuous, so $\mathcal E_A$ is a topological
relation-sufficient extension.
\end{proposition}

\begin{proof}
By the characteristic property of the quotient topology, a map out of $E_A$ is continuous iff its
composite with $e_A$ is continuous. Now $u_A\circ e_A=p$ and $\rho_\alpha\circ e_A=r_\alpha$ by
\eqref{eq:factorization}, and $p$ and $r_\alpha$ are continuous by hypothesis. Finally $e_A$ is a
quotient map by construction.
\end{proof}

\begin{theorem}[Topological universal property]
\label{thm:top-terminal}
$\mathcal E_A$, with the quotient topology, is terminal in the category of topological
relation-sufficient extensions: for every such $\mathcal E$, the unique set map
$\phi_{\mathcal E}\colon E\to E_A$ of \cref{thm:terminal} is continuous, and it is the unique
continuous morphism $\mathcal E\to\mathcal E_A$. It is a continuous surjection.
\end{theorem}

\begin{proof}
\Cref{thm:terminal} supplies the unique \emph{function} $\phi_{\mathcal E}\colon E\to E_A$ satisfying
$e_A=\phi_{\mathcal E}\circ e$, and asserts its surjectivity. For continuity: the composite
$\phi_{\mathcal E}\circ e=e_A$ is continuous (it is the quotient map onto $E_A$), and $e$ is a
quotient map by \cref{def:top-EA}; by the characteristic property of quotient maps, a function out of
$E$ whose composite with the quotient map $e$ is continuous is itself continuous. Hence
$\phi_{\mathcal E}$ is continuous. Uniqueness among \emph{continuous} morphisms is inherited from
uniqueness among all set maps, proved in \cref{thm:terminal}\,(v). Finally $\mathcal E_A$ is itself an
object of the category by \cref{prop:top-structure}.
\end{proof}

\subsection{Quotient topology versus subspace topology}

$E_A$ carries two a priori distinct topologies: the quotient topology of \cref{def:top-EA}, and the
subspace topology inherited from $T\times V_A$ via the canonical injection
\begin{equation}
\label{eq:iota}
\iota\colon E_A\hookrightarrow T\times V_A .
\end{equation}
These must not be conflated. The universal property of \cref{thm:top-terminal} is a statement about
the \emph{quotient} topology, and it is the quotient topology that the completion carries by
definition.

\begin{proposition}[Comparison, and the embedding criterion]
\label{prop:embedding}
Equip $E_A$ with the quotient topology.
\begin{enumerate}
\item\label{it:iota-cts} The injection $\iota$ of \eqref{eq:iota} is continuous. Consequently the
quotient topology is \emph{finer than or equal to} the subspace topology on $\Image(j_A)$.
\item\label{it:iota-embed} $\iota$ is a topological embedding if and only if the quotient topology
coincides with the subspace topology; equivalently, if and only if the corestriction
$j_A\colon S\to\Image(j_A)$ is a quotient map when $\Image(j_A)$ carries the product subspace
topology.
\end{enumerate}
\end{proposition}

\begin{proof}
\labelcref{it:iota-cts} We have $\iota\circ e_A=j_A$, which is continuous; since $e_A$ is a quotient
map, $\iota$ is continuous by the characteristic property. Continuity of $\iota$ says exactly that
every subspace-open subset of $E_A$ (i.e.\ every $\iota^{-1}(W)$ with $W$ open in $T\times V_A$) is
quotient-open; hence the quotient topology contains the subspace topology.

\labelcref{it:iota-embed} $\iota$ is injective and continuous, so it is an embedding iff it is open
onto its image, i.e.\ iff every quotient-open subset of $E_A$ is subspace-open. Combined with
\labelcref{it:iota-cts}, this holds iff the two topologies are equal.

For the reformulation, write $I:=\Image(j_A)$ with the subspace topology and let
$\bar j_A\colon S\to I$ be the corestriction of $j_A$, so that as \emph{functions}
$\bar j_A=\iota\circ e_A$ with $\iota$ now regarded as the identity of underlying sets $E_A\to I$.
Suppose the topologies agree. Then $\iota$ is a homeomorphism and $e_A$ is a quotient map, so
$\bar j_A$ is a quotient map. Conversely, suppose $\bar j_A$ is a quotient map. A subset $U\subseteq
E_A$ is quotient-open iff $e_A^{-1}(U)$ is open in $S$; and $e_A^{-1}(U)=\bar j_A^{-1}(U)$ as subsets
of $S$, since $e_A$ and $\bar j_A$ agree as functions into the common underlying set. Because
$\bar j_A$ is a quotient map, $\bar j_A^{-1}(U)$ open in $S$ implies $U$ open in $I$, i.e.\
subspace-open. Hence quotient-open implies subspace-open, and with \labelcref{it:iota-cts} the two
topologies coincide.
\end{proof}

\begin{corollary}[Compact--Hausdorff sufficient condition]
\label{cor:compact-hausdorff}
Suppose $S$ is compact and $T\times V_A$ is Hausdorff. Then the quotient and subspace topologies on
$E_A$ agree, and $\iota\colon E_A\hookrightarrow T\times V_A$ is a topological embedding; in
particular $E_A$ is a compact Hausdorff space.
\end{corollary}

\begin{proof}
Write $I=\Image(j_A)\subseteq T\times V_A$ with the subspace topology, and let
$\bar j_A\colon S\to I$ be the (continuous, surjective) corestriction of $j_A$. As a subspace of a
Hausdorff space, $I$ is Hausdorff. We show $\bar j_A$ is a closed map, hence a quotient map, and then
apply \cref{prop:embedding}\,\labelcref{it:iota-embed}.

Let $C\subseteq S$ be closed. Since $S$ is compact, $C$ is compact; since $\bar j_A$ is continuous,
$\bar j_A(C)$ is a compact subset of $I$; since $I$ is Hausdorff, the compact subset $\bar j_A(C)$ is
closed in $I$. Thus $\bar j_A$ is a continuous closed surjection, and every continuous closed
surjection is a quotient map: if $\bar j_A^{-1}(U)$ is open then $S\setminus\bar j_A^{-1}(U)$ is
closed, its image $\bar j_A(S\setminus\bar j_A^{-1}(U))=I\setminus U$ is closed by surjectivity, so
$U$ is open. By \cref{prop:embedding}, the topologies agree and $\iota$ is an embedding. Finally
$E_A\cong I$ is compact (continuous image of a compact space) and Hausdorff (subspace of a Hausdorff
space).
\end{proof}

\begin{remark}
\label{rem:no-agreement-without-hypotheses}
Without hypotheses of this kind the two topologies may genuinely differ, and the abstract completion
does not require them to agree: \cref{thm:top-terminal} holds for the quotient topology regardless.
One should not assert the agreement, nor treat $E_A$ as a subspace of $T\times V_A$ in the
topological sense, without verifying \cref{prop:embedding}\,\labelcref{it:iota-embed} or a sufficient
condition such as \cref{cor:compact-hausdorff}.
\end{remark}

\section{Worked example: flat \texorpdfstring{$U(1)$}{U(1)} connections on the circle}
\label{sec:u1}

We now exhibit every phenomenon of the general theory in a single, completely explicit gauge-theoretic
example. It is the extreme case of readout loss: the readout is \emph{constant}, so it distinguishes
nothing at all, and the completion must supply everything.

\subsection{The reduced source and the readout}

Consider $U(1)$ connections on the trivial bundle over the circle $S^1$, restricted to the flat ones.
Gauge reduction \eqref{eq:gauge-quotient} identifies gauge-equivalent connections; on the circle, the
resulting reduced moduli space of flat connections is classified by the holonomy around the loop, so
that
\begin{equation}
\label{eq:Sflat}
S_{\mathrm{flat}}\ \cong\ U(1),
\qquad
z=e^{i\theta},\quad\theta\in\mathbb{R}/2\pi\mathbb{Z} .
\end{equation}
This is the standard classification of flat connections on the circle by their holonomy
(see, e.g., \cite{KobayashiNomizu,WuYang}); we take it as given, and stress that $S_{\mathrm{flat}}$
is already the \emph{reduced} source in the sense of \cref{sec:setup}: gauge potentials and gauge
transformations have been quotiented away, and it is not available to us to reintroduce them.

The readout is the curvature. For flat connections the curvature vanishes identically, so
\begin{equation}
\label{eq:pF}
p_F\colon S_{\mathrm{flat}}\longrightarrow T=\{0\},
\qquad
p_F(z)=0 .
\end{equation}
The readout fiber is the whole source: $F_0=S_{\mathrm{flat}}\cong U(1)$. Every nontrivial
distinction on $S_{\mathrm{flat}}$ therefore fails to descend along $p_F$, and $\sim_{p_F}$ is the
total relation. Note carefully what is \emph{not} being claimed: we do not claim that curvature is
the complete physical observable content of a gauge field. On the contrary, the example exists
precisely because it is not---the Aharonov--Bohm effect \cite{AharonovBohm} and the Wu--Yang analysis
\cite{WuYang} are the physical statement that holonomy carries information the field strength does
not. The readout $p_F$ is a readout, not an ontology.

\subsection{Transport completion}

Nominate the single relation
\begin{equation}
\label{eq:rhol}
r_{\mathrm{hol}}\colon S_{\mathrm{flat}}\to U(1),
\qquad
r_{\mathrm{hol}}(z)=z ,
\end{equation}
the oriented holonomy itself. It is eligible in the sense of \cref{def:eligible}: holonomy around a
fixed loop is gauge-invariant (for the abelian group $U(1)$ it is invariant outright, not merely up
to conjugacy), so it is a function on the reduced source. It is nominated because oriented parallel
transport is what enters targets involving directed transport of charge around the loop.

By \cref{def:EA}, the completion is
\begin{equation}
\label{eq:Ehol}
E_{\mathrm{hol}}
=\Image\bigl(p_F,\ r_{\mathrm{hol}}\bigr)
=\{(0,z):z\in U(1)\}
\ \cong\ U(1),
\end{equation}
and $e_{\mathrm{hol}}\colon S_{\mathrm{flat}}\to E_{\mathrm{hol}}$ is a bijection. Here $S$ is compact
and $T\times V=\{0\}\times U(1)$ is Hausdorff, so \cref{cor:compact-hausdorff} applies: the quotient
and subspace topologies agree and \eqref{eq:Ehol} is a homeomorphism. By
\cref{prop:reconstruction} below, the nomination $\{p_F,r_{\mathrm{hol}}\}$ separates points of
$S_{\mathrm{flat}}$, so this is a case in which completion happens to reconstruct the whole reduced
source. It does \emph{not} restore gauge potentials or gauge representatives---by \cref{cor:gauge} it
cannot---and it restores nothing beyond the reduced source.

\subsection{Spectral completion}

Now nominate instead the charged-scalar spectral target of Paper~2 \cite{Paper2}. For a scalar of
unit charge on $S^1$ minimally coupled to the flat connection with holonomy $e^{i\theta}$, let
$H_\theta$ denote the associated Laplace-type operator. The result we import from Paper~2 is
\begin{equation}
\label{eq:spec}
\Spec(H_\theta)=\Spec(H_{\theta'})
\qquad\Longleftrightarrow\qquad
\theta'\equiv\pm\theta \pmod{2\pi},
\end{equation}
with multiplicities. We do not reprove \eqref{eq:spec} here; it is exactly the Paper~2 input, and
the new content of the present section is the family-level completion and refinement structure built
on top of it.

Take $\mathcal Q=\{Q_{\mathrm{spec}}\}$ with answer map
$a_{\mathrm{spec}}(z)=\Spec(H_\theta)$ for $z=e^{i\theta}$ (a well-defined function of $z$, since
$\theta\mapsto H_\theta$ depends on $\theta$ only mod $2\pi$). By \eqref{eq:spec}, the induced
equivalence relation on $S_{\mathrm{flat}}$ is
\begin{equation}
\label{eq:inversion}
z\ \sim_{\mathrm{spec}}\ z'
\qquad\Longleftrightarrow\qquad
z'=z\ \text{ or }\ z'=z^{-1},
\end{equation}
the readout clause $p_F(z)=p_F(z')$ being vacuous. Hence, by \cref{prop:target-completion},
\begin{equation}
\label{eq:Espec-quotient}
E_{\mathrm{spec}}\ \cong\ S_{\mathrm{flat}}/\!\sim_{\mathrm{spec}}\ =\ U(1)/(z\sim z^{-1}).
\end{equation}

\begin{proposition}[The spectral completion is an interval]
\label{prop:interval}
Give $U(1)/(z\sim z^{-1})$ the quotient topology. Then the assignment
\begin{equation}
\label{eq:cos}
c\colon U(1)/(z\sim z^{-1})\longrightarrow[-1,1],
\qquad
c\bigl([z]\bigr)=\Real z=\cos\theta
\qquad(z=e^{i\theta}),
\end{equation}
is a well-defined homeomorphism. Hence $E_{\mathrm{spec}}\cong[-1,1]$.
\end{proposition}

\begin{proof}
\emph{Well defined.} $\Real(z^{-1})=\Real(\bar z)=\Real z$ for $z\in U(1)$, so the formula does not
depend on the representative of the class.

\emph{Injective.} Suppose $\cos\theta=\cos\theta'$ with $\theta,\theta'\in\mathbb{R}/2\pi\mathbb{Z}$.
Then $\theta'\equiv\pm\theta\pmod{2\pi}$, i.e.\ $e^{i\theta'}=e^{\pm i\theta}$, i.e.\ $z'\in\{z,z^{-1}\}$,
i.e.\ $[z']=[z]$.

\emph{Surjective.} For $x\in[-1,1]$ set $\theta=\arccos x\in[0,\pi]$; then $c([e^{i\theta}])=x$.

\emph{Continuous.} The composite of $c$ with the quotient map $U(1)\to U(1)/(z\sim z^{-1})$ is
$z\mapsto\Real z$, which is continuous. By the characteristic property of the quotient topology, $c$
is continuous.

\emph{Homeomorphism.} $U(1)$ is compact, hence so is its continuous image $U(1)/(z\sim z^{-1})$; and
$[-1,1]$ is Hausdorff. A continuous bijection from a compact space to a Hausdorff space is a
homeomorphism: it is a closed map, since closed subsets of a compact space are compact, their
continuous images are compact, and compact subsets of a Hausdorff space are closed; a closed
continuous bijection is a homeomorphism.
\end{proof}

\subsection{The refinement map, concretely}

The relation family $\{r_{\mathrm{hol}}\}$ mediates the spectral target in the sense of
\cref{def:mediation}: by \eqref{eq:spec} the spectrum is a function of the holonomy, so there is
$\Phi_{\mathrm{spec}}\colon T\times U(1)\to A_{Q_{\mathrm{spec}}}$ with
$a_{\mathrm{spec}}(z)=\Phi_{\mathrm{spec}}(p_F(z),r_{\mathrm{hol}}(z))$. \Cref{thm:chi} therefore
supplies a unique surjection $\chi\colon E_{\mathrm{hol}}\to E_{\mathrm{spec}}$, and under the
identifications \eqref{eq:Ehol} and \cref{prop:interval} it is exactly
\begin{equation}
\label{eq:chi-u1}
\chi\colon U(1)\longrightarrow[-1,1],
\qquad
e^{i\theta}\longmapsto\cos\theta .
\end{equation}
It is surjective, and it is \emph{not} injective: $\chi(e^{i\theta})=\chi(e^{-i\theta})$ while
$e^{i\theta}\neq e^{-i\theta}$ whenever $\theta\not\equiv0,\pi$. So this is a concrete realization of
the strict inequality permitted by \eqref{eq:chi-order}: here $\sim_A\ \subsetneq\ \sim_{\mathcal Q}$,
and the target completion is strictly coarser than the relation completion.

The physics of the strictness is transparent. The full oriented holonomy retains an
\emph{orientation}: it distinguishes transport of charge around the loop in one sense from transport
in the other, i.e.\ $z$ from $z^{-1}$. The charged-scalar spectrum does not require that
information---by \eqref{eq:spec} it is invariant under $\theta\mapsto-\theta$---and so the minimal
completion \emph{for that target} discards it. Each completion is minimal for its own nomination.
Which is appropriate depends on which targets one has nominated, and that is a physical decision made
before the mathematics begins, not one the mathematics can make.

\begin{remark}
\label{rem:u1-endpoints}
The two fixed points of the inversion, $z=1$ ($\theta=0$) and $z=-1$ ($\theta=\pi$), are exactly the
points where $\chi$ fails to be two-to-one; they map to the endpoints $\pm1$ of the interval. They are
the connections whose holonomy is its own inverse, i.e.\ for which orientation carries no information.
\end{remark}

\section{The exact scope of reconstruction}
\label{sec:reconstruction}

Completion always reconstructs the joint readout/relation signature; that is what it is. Whether it
reconstructs the source state is a separate question with a sharp answer.

\begin{proposition}[Reconstruction criterion]
\label{prop:reconstruction}
The completion map $e_A\colon S\to E_A$ is injective---equivalently, $\sim_A\ =\ \Delta_S$;
equivalently, $e_A$ is a bijection onto $E_A$---if and only if the family
\begin{equation}
\label{eq:separating}
\{p\}\cup\{r_\alpha:\alpha\in A\}
\end{equation}
\emph{separates the points of $S$}: for all $s\neq s'$ in $S$, either $p(s)\neq p(s')$ or
$r_\alpha(s)\neq r_\alpha(s')$ for some $\alpha\in A$.
\end{proposition}

\begin{proof}
By \cref{def:EA} and \eqref{eq:simA}, for all $s,s'\in S$,
\begin{align*}
e_A(s)=e_A(s')
&\iff j_A(s)=j_A(s')\\
&\iff \bigl[\,p(s)=p(s')\ \text{ and }\ r_\alpha(s)=r_\alpha(s')\ \ \forall\alpha\in A\,\bigr]\\
&\iff s\sim_A s' .
\end{align*}
Hence $e_A$ is injective iff ($s\sim_A s'\Rightarrow s=s'$), iff $\sim_A\ =\ \Delta_S$ (the reverse
inclusion always holding by reflexivity), iff no pair of distinct points is identified by the whole
family---which is the stated separation property, in contrapositive form. Surjectivity of $e_A$ is
automatic (\cref{prop:factorization}), so injectivity is equivalent to bijectivity.
\end{proof}

\begin{remark}[Two senses of ``reconstruction''; only one is automatic]
\label{rem:two-reconstructions}
Distinguish carefully:
\begin{enumerate}
\item \emph{Reconstruction of the joint signature.} The completed state determines the readout value
and every nominated relation value: this is \eqref{eq:factorization}, and it holds \emph{always}, by
construction. It is the entire purpose of the object.
\item \emph{Reconstruction of the reduced source state.} The completed state determines $s\in S$
itself. This holds \emph{if and only if} the nominated family separates points, by
\cref{prop:reconstruction}, and in general it fails---necessarily so, since by \cref{cor:coarsest}
$E_A$ is the coarsest sufficient extension and therefore retains nothing the nomination does not
force it to retain.
\end{enumerate}
Neither sense licenses a third: recovery of \emph{raw} configuration data in $\Omega$, or selection of
a preferred representative of a physical equivalence class or of a readout fiber. \Cref{cor:gauge}
forbids the former and \cref{rem:no-selector} the latter. The word ``reconstruction'' should not be
used in this subject without saying which of (1) and (2) is meant.
\end{remark}

\Cref{sec:u1} contains an instance of each case: the transport nomination $\{p_F,r_{\mathrm{hol}}\}$
separates points of $S_{\mathrm{flat}}$, so $e_{\mathrm{hol}}$ is bijective and completion recovers
the whole reduced source; the spectral nomination does not separate $z$ from $z^{-1}$, so the
spectral completion is a proper quotient, and no amount of further analysis of the spectrum will
recover the orientation, because the orientation is not there to be recovered.

\section{Discussion and scope}
\label{sec:discussion}

\subsection{What has been established}

Given a reduced source space $S$, a surjective readout $p\colon S\to T$, and a set-indexed family of
eligible source relations $\{r_\alpha\colon S\to V_\alpha\}_{\alpha\in A}$, there exists a completion
$E_A=\Image(p,R_A)$ of the readout on which every nominated relation becomes well defined; it is
canonically the quotient $S/\!\sim_A$; it is terminal among surjective relation-sufficient extensions
and therefore unique up to unique isomorphism; it is the \emph{coarsest} such extension, in the exact
sense $\sim_e\ \subseteq\ \sim_A$ for every sufficient $e$; its fibers over $T$ are the realized
relation signatures; it is functorial and cocyclic in the nomination; it is insensitive to faithful
relabelling of relation values; it admits a canonical surjection onto the completion determined by
any family of targets that the relations mediate; it restores no gauge redundancy; and in the
topological setting it is terminal for the quotient topology, embedding into $T\times V_A$ exactly
when the quotient and subspace topologies agree.

Together with \cite{Paper1} (obstruction) and \cite{Paper2} (target-relative necessity), this
supplies the constructive step of the sequence: the object whose existence, uniqueness, and
minimality were previously only desiderata.

\subsection{The three standing confusions}

\emph{Minimality is not smallness.} \Cref{rem:minimality}: minimality is least retained distinction,
an order-theoretic statement about equivalence relations, not a cardinality bound and not an
optimality claim in any other metric.

\emph{Completion is not reconstruction.} \Cref{rem:two-reconstructions}: the completed state always
determines the nominated signature and determines the source state only under a separation hypothesis
that must be checked, not assumed.

\emph{Coordinates are relative to the nomination.} \Cref{rem:not-unique-coords} and
\cref{prop:reparam}: the nominated relations are a presentation of the quotient, and irredundancy
constrains that presentation without making it unique. Relative to a fixed nominated family, $E_A$
does have canonical structure maps $u_A$ and $(\rho_\alpha)_{\alpha\in A}$; what it lacks is a
family-independent, intrinsically unique coordinate presentation. Different faithful
reparameterizations, and different irredundant families, may present the same quotient.

\subsection{What has not been established}

The theorem is model-relative, and the following are \emph{not} consequences of it.

That the universe is a projection of a hidden source. Nothing here bears on that; $S$ is an assumed
structure, and the theorem tells one what follows from assuming it, not that it exists.

That the source space is empirically real. The construction is available whether $S$ is a physical
posit, a redundant description, or a bookkeeping device.

That every lost distinction ought to be restored. On the contrary: the entire point of
\cref{thm:chi} and of \cref{sec:u1} is that a distinction is worth restoring only relative to a
nominated target, and that nominating more than the target requires yields a completion that is
finer than necessary---correct, but not minimal for that target.

That the full source can in general be reconstructed. \Cref{prop:reconstruction} makes this a
hypothesis, and the spectral example shows it can fail.

That the nominated relation family is physically privileged. This is the load-bearing external input,
and the mathematics cannot supply it. A relation is not physical because it varies on a readout
fiber; if it were, every gauge artifact would be an observable. The nomination must be justified by
the physics of the targets, and this paper's theorems are conditional on it in exactly that one
respect---which is, we think, where the conditionality belongs.

\subsection{Outlook}

The present construction is purely set-theoretic and, in \cref{sec:topology}, topological. It says
nothing about when $E_A$ carries additional structure---smooth, symplectic, metric---or when the
relation signatures arise as genuine geometric invariants rather than as formal coordinates. The
companion manuscript in preparation on geometric realization takes up exactly that question, on
restricted spaces of metrics, bundles, connections, and matter fields. We have deliberately imported
none of that machinery here; the theorems above should stand or fall on their own, and they require
nothing beyond images, quotients, and quotient topologies.

\section{Conclusion}
\label{sec:conclusion}

Source completion is neither arbitrary enrichment of the readout nor selection of a preferred state
from each readout fiber. It is the canonical quotient of the reduced source that preserves the
original readout while separating exactly those source states that the nominated family of invariant
relations distinguishes---and no others. The joint-image construction proves existence; the terminal
universal property, with its arrow running $E\to E_A$ from every sufficient extension onto the
completion, proves minimality and uniqueness up to unique isomorphism; the indispensability criterion
identifies which nominated coordinates do genuine work; the refinement tower organizes nominations by
information content; the target-completion surjection separates ``what the relations say'' from ``what
the answers need''; the gauge boundary confirms that completion cannot undo reduction; and the
topological theorem carries the whole construction into continuous settings without conflating the
quotient and subspace topologies. What remains outside the theorem---whether a given source is real,
and whether a given nomination is physical---remains outside it deliberately.

\begin{thebibliography}{9}

\bibitem{Paper1}
J.~Rodgers,
\emph{Source--Readout Non-Equivalence: Descent and Equivariant Reconstruction Obstructions},
Zenodo, 2026.
\href{https://doi.org/10.5281/zenodo.21369967}{doi:10.5281/zenodo.21369967}.

\bibitem{Paper2}
J.~Rodgers,
\emph{Target-Relative Necessity of Completion: When Readout Loss Obstructs, and What a Sufficient
Extension Must Retain},
Zenodo, 2026.
\href{https://doi.org/10.5281/zenodo.21370501}{doi:10.5281/zenodo.21370501}.

\bibitem{MacLane}
S.~Mac~Lane,
\emph{Categories for the Working Mathematician},
2nd ed., Graduate Texts in Mathematics 5, Springer-Verlag, New York, 1998.

\bibitem{Munkres}
J.~R.~Munkres,
\emph{Topology},
2nd ed., Prentice Hall, Upper Saddle River, NJ, 2000.

\bibitem{KobayashiNomizu}
S.~Kobayashi and K.~Nomizu,
\emph{Foundations of Differential Geometry}, Vol.~I,
Interscience Publishers, New York, 1963.

\bibitem{AharonovBohm}
Y.~Aharonov and D.~Bohm,
\emph{Significance of electromagnetic potentials in the quantum theory},
Phys. Rev. \textbf{115} (1959), 485--491.

\bibitem{WuYang}
T.~T.~Wu and C.~N.~Yang,
\emph{Concept of nonintegrable phase factors and global formulation of gauge fields},
Phys. Rev. D \textbf{12} (1975), 3845--3857.

\end{thebibliography}

\end{document}
